\section{Research goals}\label{sec:goal}

The goal of this research is to establish best practices for modeling large railway systems in answer set programming.
Our initial goal has been to understand large railway systems by leveraging the virtual railway environment generator, Flatland.
In doing so, developing a comprehensive formalization of the problem, as well as initial encodings, has enabled us to see how the problem of railway management differs from those in similar domains.
Ultimately, this understanding can be translated to more general problems faced by railway operators.
Narayanaswami and Rangaraj classify problems in railway operations into broad categories from planning to execution \cite{narran2011a}.
Specific problems from across that spectrum that are worth addressing include
\begin{itemize}
	\item routing and scheduling
	\item acute replanning during adverse events
	\item valid switch position logic (such as that used by the ILTIS system)
	\item energy utilization
	\item route synchronization
\end{itemize}
Although solutions to many issues in the realm of transportation are rooted in public policy,
it is my hope that my research can contribute to the solving of problems that require or can benefit from technological solutions.

%
%%% Local Variables:
%%% mode: LaTeX
%%% TeX-master: "paper"
%%% End:
